\documentclass[a4paper]{article}
\XeTeXlinebreaklocale "ja"
\usepackage{xltxtra}

% fonts
\usepackage{fontspec}
\setmainfont[Scale=MatchLowercase]{NotoSerifJP-Regular}
\setsansfont[Scale=MatchLowercase]{NotoSansJP-Regular}
\setmonofont[Scale=MatchLowercase]{CourierPrime-Regular}

% bibtex
\usepackage{cite}

% 画像
\usepackage{graphicx}

\begin{document}

  \title{ReXeTeXeR}
  \author{Takahito Sueda}
  \maketitle

  \section{はじめに}
  {\XeTeX} でいい感じにするやつです。

  \section{数式}
  普通に{\TeX}で書けます。
  \begin{eqnarray}
    2x_1 + x_2 & = & 5 \\
    2x_2 & = & 2
  \end{eqnarray}

  \section{引用}
  bibtexを使えます。引用箇所は↓の感じです。Referencesが最後のとこにあります。\\
  引用テストDL\cite{lecun2015deep} \\
  引用テストML\cite{michie1994machine}

  \section{画像}
  graphicx使っていけます。
  \begin{center}
    \includegraphics[width=10cm]{img/logo.png} \\
    ReXeTeXeRのぶちかっこいいロゴ
  \end{center}

  \section{自動監視テスト}
  なんか書くと、コンテナ内のシェルスクリプトがこのファイルの変更を検知して、texのコンパイルコマンドが走ります。

  % bibtex
  \bibliographystyle{junsrt}
  \bibliography{ref.bib}
\end{document}
